%
\chapter{Problem Formulation}
%

The optimization of charging infrastructure placement presents a complex multi-objective challenge requiring careful mathematical formulation. We developed a \glsxtrfull{mip} that balances coverage maximization with cost efficiency while adhering to practical infrastructure constraints. This formulation addresses the dual challenges of expanding charging coverage while maintaining cost efficiency through strategic station placement and infrastructure upgrades.

\section{Data}

% \subsection{Sets and Indices}

% The existing infrastructure is divided into Level 3 stations ($S^3$), and Level 2 stations ($S^2$) which are eligible for being upgraded to Level 3 charging capabilities.

% The foundation of the model lies in its geographical representation, where potential locations for new stations ($P$) and population-based demand points ($D$) are defined. The normalized population weight, $w_j$ reflects the relative importance of each demand point based on demographic factors. For coverage assessment, we define $P^2_j$ and $P^3_j$ as the sets of potential locations that can serve demand points $j, j \in D$ within the coverage radii of Level 2 and Level 3 charging ports, respectively. Each demand point carries a normalized weight ($w_j$).

% \subsection{Coverage Parameters}

% The coverage parameters define how charging stations serve demand points. Coverage radii $r_2$ and $r_3$ determine the service area of Level 2 and Level 3 stations. Minimum coverage requirements $\alpha_2$ and $\alpha_3$ ensure adequate service levels across the region. These values were set in our base model (\verb|configs\base.json|) based on typical urban travel patterns.

% For a tabulated list of coverage parameters, see Table \ref{tab:coverage-parameters} in the Appendix.

% \subsection{Cost Parameters}

% Cost parameters capture the financial aspects of infrastructure development. Fixed costs $c_2^s$ and $c_3^s$ represent the base installation costs for new stations, while per-port costs $c_2^p$ and $c_3^p$ account for individual charging unit installations. The resale value factor $\beta$ determines the recoverable value when upgrading Level 2 stations. All expenditures must remain within the total available budget $B$. These values were set in our base model (\verb|configs\base.json|) based on typical figures from \glsxtrshort{ev} charger articles. 

% For a detailed tabulated list of cost parameters, see Table \ref{tab:cost-parameters} in the Appendix \cite{a2021_kitchenerwaterloocambridge} \cite{a2024_kitchenerwaterloocambridgeguelph} \cite{a2021_kitchenerwaterloocambridgeguelph} \cite{dinello_2022_what} \cite{a2024_ev} \cite{charging_2023_cost}.

% \subsection{Infrastructure Parameters}

% Infrastructure parameters define the technical specifications and limitations. Minimum port requirements $p_2$ and $p_3$ ensure sufficient charging capacity at each station. Power requirements $g_2$ and $g_3$ specify the electrical load per charging port, which must not exceed the local grid capacity $G_i$ at any location.

% For a detailed tabulated list of infrastructure parameters, see Table \ref{tab:infrastructure-parameters} in the Appendix.

\begin{enumerate}

    \item \textbf{Sets and Indices:} Our model incorporates existing infrastructure categorization into \gls{l3} stations ($S^3$) and \gls{l2} stations ($S^2$) eligible for upgrades. The geographical framework encompasses potential locations for new stations ($P$) and population-based demand points ($D$), with each demand point carrying a normalized weight ($w_j$) reflecting its demographic significance. Coverage assessment utilizes sets $P^2_j$ and $P^3_j$, representing potential locations capable of serving demand points within \gls{l2} and \gls{l3} charging radii, respectively. 
    
    For a detailed listing of sets and indices used in our formulation, see Table \ref{tab:sets-and-indices} in the Appendix.

    \item \textbf{Coverage parameters} define service areas through radii $r^2$ and $r^3$ for \gls{l2} and \gls{l3} stations, with minimum coverage requirements $\alpha^2$ and $\alpha^3$ ensuring adequate service distribution. 
    
    These values, detailed in Table \ref{tab:coverage-parameters} of the Appendix, were established based on standard urban mobility patterns and charging technology specifications from SparkCharge (2024). \cite{spark2024}

    \item \textbf{Cost parameters}, outlined in Table \ref{tab:cost-parameters} of the Appendix, encompass fixed costs $c_2^s$ and $c_3^s$ for new station installations, alongside per-port costs $c_2^p$ and $c_3^p$. The model incorporates a resale value factor $\beta$ for equipment recovery during \gls{l2} to \gls{l3} upgrades, with all expenditures constrained by the total available budget $B$. \cite{joyce2023}
    
    These financial parameters were derived from current market data and industry cost analyses. \cite{future2022}

    \item \textbf{Infrastructure parameters} define technical specifications, including level-specific minimum charging port requirements $p_2$ and $p_3$, and power requirements $g_2$ and $g_3$ per charging port. These specifications must respect local grid capacity $G_i$ at each location, as detailed in Table \ref{tab:infrastructure-parameters} of the Appendix. 
    
    The values align with grid capacity assessments from the IESO Kitchener-Waterloo-Cambridge-Guelph Region Integrated Regional Resource Plan \cite{ieso2021}.
\end{enumerate}

\section{Decision Variables}

% The decision-making framework employs binary variables to determine infrastructure changes: $x^2_i$ for new Level 2 stations, $x^3_i$ for new Level 3 stations, and $u_k$ to upgrade existing Level 2 stations to Level 3. Coverage achievement is tracked through $y^2_j$ and $y^3_j$, indicating whether demand point $j$ receives Level 2 or Level 3 service.

% For a detailed tabulated list of decision variables, see Table \ref{tab:decision-variables} in the Appendix.

Our model employs binary decision variables to determine infrastructure modifications. Variables $x^2_i$ and $x^3_i$ govern new \gls{l2} and \gls{l3} station placements, respectively, while $u_k$ determines upgrades of existing \gls{l2} stations to \gls{l3} capability. Coverage achievement tracking utilizes variables $y^2_j$ and $y^3_j$ to indicate service provision at each demand point for \gls{l2} and \gls{l3} charging facilities (stations and/or ports), respectively. 

For a comprehensive listing of decision variables and their constraints, refer to Table \ref{tab:decision-variables} in the Appendix.

\section{Constraints}

% Our model incorporates four essential constraint categories that ensure practical feasibility and operational efficiency.

% \subsection{Station Placement Type}

% To prevent redundant infrastructure and optimize resource allocation, we ensure that at most one type of charging station (Level 2 or Level 3) can be installed at any potential location. See Inequality \ref{eq:const-type}.

% \subsection{Budget}

% The total expenditure $C$ must remain within the available budget $B$. This includes installation costs for new stations and charging ports, and the net cost of upgrades after accounting for the resale value of existing Level 2 equipment. See Inequality \ref{eq:const-budget}.

% \subsection{Minimum Coverage Targets}

% Our model strives to ensure adequate service coverage across the region. For Level 2 coverage, a demand point is served if it's within range of a Level 2 charging port. This could be satisfied by either an existing new Level 2 station, a new Level 2 station, or a Level 2 station that was upgraded to a Level 3 station through port retention. See Inequality \ref{eq:const-coverage-l2}.

% Similarly, for Level 3 coverage, a demand point is served if it is within the range of a Level 3 charging port. This could be satisfied by either an existing Level 3 charging station, a new Level 3 charging station, or an upgraded Level 2 station within the specified radius. See Inequality \ref{eq:const-coverage-l3}.

% To ensure equitable access, we require minimum coverage levels weighted by demographic importance. See Inequalies \ref{eq:const-coverage-req-l2} and \ref{eq:const-coverage-req-l3}.

% \subsection{Grid / Power Capacity}

% The power requirements of charging stations must not exceed the local electrical grid capacity. This accounts for both the base load of Level 2 ports and the additional demand from Level 3 charging. See Inequality \ref{eq:const-grid}.

Four essential constraint categories ensure practical feasibility and operational efficiency in our model. 

\begin{enumerate}
    \item The \textbf{Station Placement Type constraints} prevent redundant infrastructure by limiting each location to at most one charging station type. See Inequality \ref{eq:const-type}.

    \begin{align}
        x^2_i + x^3_i \leq 1, \quad \forall i \in P
        \label{eq:const-type}
    \end{align}

    \item The \textbf{Budget constraint} encompasses installation costs for new stations and charging ports, alongside the net cost of upgrades after accounting for equipment resale value. See Inequality \ref{eq:const-budget}.

    \begin{align}
        C =
        & \sum_{i}(c_2^s \cdot x^2_i + p_2 \cdot c_2^p \cdot x^2_i) + \nonumber \label{eq:net-cost} \\ 
        & \sum_{i}(c_3^s \cdot x^3_i + p_3 \cdot c_3^p \cdot x^3_i) + \nonumber \\
        & \sum_{i \in S^2}(u_i (c_3^s + p_3 \cdot c_3^p - \beta \cdot c_2^s - \beta \cdot p_i \cdot c_2^p)) \\
        \nonumber \\
        C \leq &  B \label{eq:const-budget}
    \end{align}

    \item \textbf{Coverage constraints} ensure adequate service levels across the region. \gls{l2} coverage requirements consider both new installations and retained ports from upgraded stations, while \gls{l3} coverage accounts for new installations and upgraded facilities. See Inequalities \ref{eq:const-coverage-l2} and \ref{eq:const-coverage-l3}.

    \begin{align}
        y^2_j & \leq \sum_{i \in P^2_j} x^2_i + \sum_{i \in S^2} (1-u_i) & \quad \forall j \in D 
        \label{eq:const-coverage-l2} \\
        y^3_j & \leq \sum_{i \in P^3_j} x^3_i + \sum_{i \in S^2} u_i & \quad \forall j \in D 
        \label{eq:const-coverage-l3}
    \end{align}
    
    These constraints utilize population-weighted requirements to ensure equitable service distribution. See Inequalities \ref{eq:const-coverage-req-l2} and \ref{eq:const-coverage-req-l3}.

    \begin{align}
        \sum_{j \in D}(w_j \cdot y^2_j) & \geq \alpha_2
        \label {eq:const-coverage-req-l2} \\
        \sum_{j \in D}(w_j \cdot y^3_j) & \geq \alpha_3
        \label{eq:const-coverage-req-l3}
    \end{align}

    \item \textbf{Grid capacity constraints} ensure that power requirements at each location remain within local electrical infrastructure limitations. These constraints consider both the base load of \gls{l2} ports and the additional demand from \gls{l3} charging capabilities, adhering to specifications from Hydro One's regional planning guidelines \cite{hydro2021} \cite{hydro2024}. See Inequality \ref{eq:const-grid}.

    \begin{equation}
        g_2 \cdot p_2 \cdot x^2_i + g_3 \cdot p_3 \cdot (x^3_i + u_i) \leq G_i, \quad \forall i \in P
        \label{eq:const-grid}
    \end{equation}
\end{enumerate} 

\section{Objective Function}

Our multi-objective function balances competing priorities through a weighted approach combining maximization of \gls{l2} service coverage (from \ref{eq:const-coverage-req-l2}), maximization of \gls{l3} service coverage (from \ref{eq:const-coverage-req-l3}), and minimization of total cost (from \ref{eq:net-cost}).

\begin{align}
    \left\{
    \begin{array}{ll}
        \max & \sum\limits_{j \in D} w_j \cdot y^2_j \\[1.5em]
        \max & \sum\limits_{j \in D} w_j \cdot y^3_j \\[1.5em]
        \min & C
    \end{array}
    \right\} \tag{3-OBJ} \label{3-OBJ}
\end{align}

\eqref{3-OBJ} is converted into a weighted linear function. The positive objective weights ($\lambda_1$, $\lambda_2$, $\lambda_3$) reflect strategic priorities while summing to unity. This formulation allows for scenario analysis through weight adjustment while maintaining model linearity.

\begin{align}
    \max \quad & \lambda_1\sum_{j \in D}(w_j \cdot y^2_j) + \lambda_2\sum_{j \in D}(w_j \cdot y^3_j) - \lambda_3C \tag{WT-LP} \label{WT-LP}
\end{align}

The weights set for our base model are outlined in Table \ref{tab:obj-wt} in the Appendix. The polarity between the coverage maximization weights and the cost minimization weights is based on the magnitude of values.

% The resulting optimization framework provides a systematic approach to the development of electric vehicle charging infrastructure, balancing immediate coverage needs with long-term sustainability while adhering to practical constraints.